\section{Source Control and the Assembly Algorithm}

\subsection{The four controlling layers}

\begin{longtable}{>{\bfseries\raggedright\arraybackslash}p{0.18\linewidth}>{\raggedright\arraybackslash}p{0.27\linewidth}>{\raggedright\arraybackslash}p{0.47\linewidth}}
\toprule
Layer & Principal source & What it decides\\
\midrule
\endhead
Calendar & Universal calendar plus the applicable diocesan, local, church, and religious calendars & Which temporal, sanctoral, proper, vigil, octave, and Saturday observances are candidates.\\
Precedence & General Rubrics 4--114, especially rubric 91 & Which candidate supplies the Office and Mass of the day, and whether another is transferred, reposed, commemorated, or omitted.\\
Mass category & General Rubrics of the Missal 269--423 and a particular formulary rubric & Whether the celebration uses the Mass of the Office, a festive, votive, Requiem, ritual, or externally solemn Mass.\\
Assembly & General Rubrics 424--516; \latin{Ritus servandus}; \latin{Ordo Missae}; \latin{Canon Missae}; selected formulary & Which texts are said, in what order, with what seasonal forms, prayers, Preface, Canon insert, dismissal, and Last Gospel.\\
\bottomrule
\end{longtable}

A particular rubric can qualify a general rule. The Nuptial Mass, Ember Saturday, Palm Sunday, Paschal Vigil, and Requiem formularies cannot be reconstructed safely from a generic ten-proper template. Read every heading and red direction attached to the selected formulary before declaring the assembly complete.

\subsection{Input sheet: facts that cannot be guessed}

Before doing any comparison, record:

\begin{enumerate}[leftmargin=1.45em]
\item civil date and liturgical season;
\item country, diocese, city or place, and the particular church or oratory;
\item title and dedication status of the church, and principal patrons;
\item whether a religious institute's calendar applies;
\item celebrant's calendar when the place rubric permits a choice;
\item whether the Mass is conventual, nonconventual sung, or Low;
\item any funeral, sacrament, blessing, procession, exposition, public devotion, external solemnity, or competent indult;
\item number of Masses in the church when a permission is number-limited.
\end{enumerate}

Rubrics 274--284 determine which calendar governs. In a public church or public oratory, the church's calendar normally binds any visiting priest. Secondary semipublic oratories and private oratories have specified alternatives. A universal calendar alone cannot disclose a parish's I-class title, the anniversary of its dedication, a diocesan patron, or a religious founder.

\subsection{Ten-step algorithm}

\begin{massbox}{Step 1: list every candidate day}
Enter the temporal day, every universal feast, every sourced particular feast, any vigil, the applicable octave day, and Our Lady on Saturday if rubric 78 applies. Do not yet choose a Mass.
\end{massbox}

\begin{massbox}{Step 2: classify the candidates}
For each candidate record its type---Sunday, feria, vigil, feast, octave day, or Saturday Office---and its own class. A Mass category has not yet been chosen, so do not label the day “a class III votive.”
\end{massbox}

\begin{massbox}{Step 3: assign precedence positions}
Locate each candidate in rubric 91. Within the same nominal class, a particular feast may precede a universal feast, a universal feast may precede an Ember feria, or a Sunday may precede a saint. Write the table position, not only the class.
\end{massbox}

\begin{massbox}{Step 4: apply express exceptions}
Test the rules for I- and II-class Sundays, feasts of the Lord, the Immaculate Conception, All Souls, Christmas octave, vigils, and the resumed Sundays after Epiphany. These rules prevent an apparently plausible result from being reached by class comparison alone.
\end{massbox}

\begin{massbox}{Step 5: dispose of the impeded item}
Use rubrics 92--102 to decide accidental transfer, perpetual reposition, commemoration, or omission. A lower position does not itself say which consequence follows.
\end{massbox}

\begin{massbox}{Step 6: calculate commemorations}
Classify each as privileged or ordinary under rubrics 106--110, observe the limits in rubric 111, the exclusions in rubric 112, and the order in rubrics 113--114. Then test whether the actual Mass is sung nonconventually, because ordinary commemorations do not enter it.
\end{massbox}

\begin{massbox}{Step 7: test the requested Mass category}
Begin with the Mass corresponding to the Office. If another Mass is requested, identify the exact rule admitting a festive, votive, Requiem, ritual, or external-solemnity Mass on this class of day. Do not reason from devotion or pastoral desirability alone.
\end{massbox}

\begin{massbox}{Step 8: select and complete the formulary}
Take the proper formulary where one exists. If it refers to a Common, supply only the missing parts from the indicated Common. Apply its seasonal alternatives and exact status as festive or votive. Keep a source line for every component.
\end{massbox}

\begin{massbox}{Step 9: merge the formulary into the Order}
Resolve the Introit doxology, Gloria, prayer set, lessons and chants, Creed, Preface, Canon inserts, Communion and Postcommunions, prayer over the people, dismissal, blessing, and Last Gospel. The complete worksheet in Section~\ref{sec:assembly-map} follows the order of celebration.
\end{massbox}

\begin{massbox}{Step 10: perform a contradiction check}
Ask whether the result uses two calendars, exceeds the prayer limit, commemorates the same Person twice, imports a Preface through a commemoration, grants a votive on a forbidden day, or carries an earlier or later edition's practice into 1962. Recheck every exceptional claim in its numbered rubric and the selected formulary.
\end{massbox}

\subsection{Why an Ordo is still necessary}

The algorithm is reproducible, but its inputs are not universally printed in one place. A properly prepared Ordo has already joined the movable temporal cycle, general calendar, national and diocesan calendar, local patronal and dedication data, and competent dispensations or indults. This manual explains how to audit that answer and how to assemble its Mass; it is not a fictional universal Ordo for every altar.
